\rok{Новембар 2024.}{I}

Први практични тест из Алгоритама и структура података

\zadatak{1}{Merge Sort}
Написати \textit{merge} методу \textit{Merge Sort} алгоритма за сортирање целобројног низа дужине \textit{n}.

\textbf{Додатне напомене за решавање задатка:}
\begin{itemize}
    \item Улаз је несортиран низ целих бројева.
    \item Излаз је сортирани низ целих бројева.
    \item У шаблону се налазе класе \texttt{RandomArrayGenerator} и \texttt{RandomNumberGen} које треба искористити за генерисање низа случајних бројева.
    \item У оквиру класе \texttt{MergeSort} написати имплементацију за методу:
    \begin{Verbatim}[fontsize=\footnotesize, numbers=left, xleftmargin=10mm]
private static void merge(long[] arr, long[] tmpArray,
                int lowerIndex, int middleIndex, int upperIndex)
    \end{Verbatim}
    \item Обезбедити код у \texttt{Main} методи за тестирање алгоритма.
\end{itemize}



\zadatak{2}{Промена редоследа n елемената стека, користећи ред}
Имплементирати функцију која за задати стек целих бројева обрће редослед само првих n елемената стека, користећи помоћни ред. У оквиру тих првих n елемената, за сваки од њих који је паран важиће да се његова вредност умањи за 2, а уколико је неки од тих првих n елемената дељив са 3 његова вредност ће се удвостручи. Уколико се након умањивања парног броја добије вредност 0, вредност тог елемента поставити на 1. Након тога, преостали елементи стека остају у истом редоследу.

\textbf{Спецификација функције:}
\begin{itemize}
    \item Функција прима број n, где је n број елемената које треба обрнути.
    \item Функција треба користити ред (\textit{queue}) као помоћну структуру.
    \item Елементи после првих n треба да остану у оригиналном редоследу.
\end{itemize}

\textbf{Додатни захтеви за решавање задатка:}
\begin{itemize}
    \item Алгоритам писати у оквиру класе \texttt{Stack}.
    \item Захтеван потпис методе:
    \begin{Verbatim}[fontsize=\footnotesize, numbers=left, xleftmargin=10mm]
public void ReverseAndModifyTopNElements(int n)
    \end{Verbatim}
\end{itemize}

\textbf{Примери:}
\begin{itemize}
    \item \textbf{Пример 1:}
    \begin{itemize}
        \item Оригинални стек: 4 3 1 7 9 4, $n = 3$
        \item Стек након обртања 3 елемента: 1 6 2 7 9 4
    \end{itemize}
    \item \textbf{Пример 2:}
    \begin{itemize}
        \item Оригинални стек: 4 2 6 10 9 4, $n = 3$
        \item Стек након обртања 3 елемента: 2 1 12 10 9 4
    \end{itemize}
    \item \textbf{Пример 3:}
    \begin{itemize}
        \item Оригинални стек: 1 2 3 4 5 6 7, $n = 8$
        \item Неважећа вредност за n.
    \end{itemize}
\end{itemize}



\zadatak{3}{Број и сума „увала" и „брегова" у једноструко повезаној листи}
У једноструко спрегнутој листи сваки од чворова може да представља увала или брег.

Чвор ће представљати \textbf{увала} уколико је његова вредност мања од оба суседна елемента (елемент пре и после посматраног чвора у оквиру листе) -- додатно, разлика у односу на оба суседа мора бити најмање 2.

Чвор ће представљати \textbf{брег} уколико је његова вредност већа од оба суседна елемента (елемент пре и после посматраног чвора у оквиру листе) -- додатно, разлика у односу на оба суседа мора бити најмање 3.

Метода враћа низ у коме прва вредност представља број увала и брегова, а друга суму њихових вредности.

Напомена: чвор може бити увала или брег само ако има претходни и следећи чвор.

\textbf{Додатни захтеви:}
\begin{itemize}
    \item Решавај овај задатак помоћу структуре листе (\textit{LinkedList}) из шаблона задатка.
    \item Имплементирај га у оквиру методе:
    \begin{Verbatim}[fontsize=\footnotesize, numbers=left, xleftmargin=10mm]
public int[] CountAndSumValleysAndPeaks()
    \end{Verbatim}
\end{itemize}

\textbf{Примери:}
\begin{itemize}
    \item \textbf{Пример 1:}
    \begin{itemize}
        \item \textbf{Улаз:} листа: $1 \rightarrow 3 \rightarrow 1 \rightarrow 4 \rightarrow 2 \rightarrow 7 \rightarrow 3 \rightarrow 6$
        \item \textbf{Излаз:} [4, 13]
        \item Објашњење: 1 је мање за 2 или више од суседних 3 и 4. 2 је мање за 2 или више од суседних 4 и 7. 3 је мање за 2 или више од суседних 7 и 6. 7 је веће за 3 или више од суседних 2 и 3. Увале: 1, 2 и 3. Брегови: 7. Укупан број је 4, а сума 1+2+3+7 је 13.
    \end{itemize}
    \item \textbf{Пример 2:}
    \begin{itemize}
        \item \textbf{Улаз:} листа: $1 \rightarrow 2 \rightarrow 2 \rightarrow 5 \rightarrow 5$
        \item \textbf{Излаз:} [-1, -1]
        \item Немамо критичне тачке, односно увале и брегове јер ниједан од чворова не задовољава постављене услове.
    \end{itemize}
\end{itemize}

\sablon{Шаблон првог теста новембар 2024. група I}{2024/Novembar/GrupaI/AISP_2025_Test1_GrupaI.linq}
